\newtcolorbox{myExBox}[1][]{
  enhanced jigsaw,
  breakable,
  parbox=false,
  attach boxed title to top text left={yshift=-3mm,yshifttext=-3mm},
  colback=ForestGreen!10,
  colframe=ForestGreen,
  colbacktitle=ForestGreen,
  #1}

\newtcolorbox{myBox}[1][]{
  enhanced jigsaw,
  lines before break=1,
  breakable,
  attach boxed title to top text left={yshift=-3mm,yshifttext=-1mm},
  colback=RoyalBlue!10,
  colframe=RoyalBlue,
  colbacktitle=RoyalBlue,
  #1}
  
\section*{Hospitationsgefäss BPA}

\begin{myExBox}[title=Beobachtungsschwerpunkte]
\textbf{Beobachtungsschwerpunkt 1}: Wie kann man an die Lebenswelt der SchülerInnen anknüfen? Wann funktioniert dies, wann weniger?

Erfahrungsgemäss ist es sehr wichtig, bei Unterrichtseinheiten gleich zu Beginn die intrinsische Motivation der SchülerInnen zu fördern, indem ein \textit{persönlicher} und \textit{konkreter} Bezug zu ihrere Lebenswelt gemacht wird. Dies scheint mir umso relevanter an Berufsschulen, wo es auch darum geht, Inhalte zu lernen, die für das Ausüben eines konkreten Berufs vonnutzen sein könnten. Im Rahmen der Hospitation möchte ich mich nicht nur darauf achten, ob solche Bezüge hergestellt wurden, sondern auch unter welchen Bedingungen sie erfolgreich zur Steigerung der Eigen-Motivation geführt haben.

\textbf{Beobachtungsschwerpunkt 2}: ...
\end{myExBox}


\begin{myExBox}[title=Persönliche Beantwortung der Fragestellungen]
Ursprüngliche Fragen:
...

Persönlichen Einschätzungen:
\begin{enumerate}
\item Ich hätte dem Unterrichtseinstieg mehr Aufmerksamkeit widmen und den SuS nicht nur dem Thema sondern auch der (vorhersehbaren) Situation besser anpassen sollen. Möglicherweise wäre es sinnvoll gewesen, wenn ich mit einer komplexen, lernprozessförderlichen Aufgabe im Sinne von \cite{tulo} begonnen hätte. Vermutlich wäre es auch sinnvoll gewesen, durch die Aufgabe, die die SuS mit dem vorhandenen Wissen noch nicht lösen können, eine ``kognitive Dissonanz'' bei den SuS, bzw. ein Gefühl der Unzufriedenheit auszulösen und somit die Lernmotivation zu steigern. Dabei würde ich in Zukunft insbesondere versuchen, einen Bezug zur Lebenswelt und dem Vorwissen der SuS herzustellen. Die Aufgabe würde dann meinen weiteren Unterricht lenken und sollte als Folge des Lernprozesses durch die SuS beantwortet, bzw. gelöst werden können. Nebst alldem sei auch meine praktische Erkenntnis erwähnt, dass ich in Zukunft versuchen werde, Hospitationen durch das Rektorat nicht mehr auf Termine direkt nach Prüfungen zu setzen.
\item Beim elementaren Programmierkenntnisse geht es aus meiner Sicht um deklaratives Wissen (Fakten) sowie prozedurales Wissen (Regeln, Handlungssequenzen). Allerdings wird das Wissen mit fortschreitendem Niveau immer komplexer und vernetzter, weshalb ich mit fortschreitendem Niveau die direkte Instruktion (z.B. neue Lerninhalte und Aufgaben präsentieren) mit der kognitiven Meisterlehre verbinden würde (Vorzeigen einer realen Anwendung, Nachahmen durch die SuS).
\item Sowohl komplexe Probleme wie auch komplexe Gestaltungsaufgaben scheinen sinnvoll, um Programmierkenntnisse zu erwerben. Ein Verlaufsmodell der direkten Instruktion würde ich wohl mit komplexen Problemen kombinieren, währenddem ich eine Meisterlehre mit komplexen Gestaltungsaufgaben kombinieren würde.
\end{enumerate}
Weiterführende Fragen:
\begin{itemize}
\item Wie kann damit umgegangen werden, dass eine Klasse aufgrund vorausgegangener Ereignisse (Anzahl Lektionen, Tageszeit, Prüfung etc.) wenig aufnahmefähig scheint?
\item Wie kann man den Unterricht noch ``retten'', wenn es anders läuft als geplant?
\item Wie kann überprüft werden, ob ein Thema verstanden wurde?
\end{itemize}

\end{myExBox}